\chapter{Propósito del manual de réplica}

Esta parte traduce las hipótesis de COBRIA a un procedimiento ejecutable en motores
NoSQL documentales. Su función no es imponer un proveedor, sino impedir que cada réplica
redefina silenciosamente la unidad de datos, el índice, el ámbito o el criterio de
equivalencia. Una comparación solo es interpretable cuando las configuraciones responden
la misma pregunta lógica y registran las diferencias físicas relevantes.

\section{Resultado esperado}

Al terminar una réplica, un tercero debe poder reconstruir cuatro productos: el conjunto
de documentos de entrada, las configuraciones comparadas, el registro de cada corrida y
el informe que enlaza cifras con observaciones originales. También debe poder comprobar
qué se excluyó, por qué se excluyó y si una enmienda ocurrió antes o después de conocer
el resultado afectado.

\section{Principios de neutralidad}

\begin{enumerate}[leftmargin=*]
  \item la semántica del dominio permanece igual entre motores;
  \item cada adaptador usa capacidades documentales normales del motor evaluado;
  \item los índices se declaran y congelan antes de medir;
  \item las proyecciones contienen únicamente los campos necesarios para la consulta;
  \item costo, frescura y reparación se informan junto con la latencia;
  \item una configuración que pierde permanece en el conjunto publicado.
\end{enumerate}

\chapter{Selección y descripción de motores documentales}

\section{Criterios de inclusión}

Un motor es elegible cuando almacena documentos u objetos agregados, permite claves y
filtros con ámbito, ofrece índices declarables, admite operaciones de actualización y
posee una vía reproducible local o administrada. La réplica debe incluir al menos dos
motores con modelos operacionales distintos para evitar que una capacidad particular se
confunda con una propiedad general de COBRIA.

\section{Ficha del motor}

Cada adaptador publica versión, modalidad de despliegue, región, consistencia observable,
límites de documento, estrategia de partición, índices, paginación, transacciones
documentales disponibles, flujos de cambios, cuotas y unidad de facturación. Si una
capacidad no existe, el protocolo registra la adaptación en vez de simular que todos los
motores son idénticos.

\begin{table}[H]\centering
\caption{Campos mínimos de la ficha del motor NoSQL.}
\begin{tabularx}{\textwidth}{p{4cm}X}\toprule
Campo & Pregunta de control\\\midrule
Identidad de versión & ¿Qué versión exacta o fecha del servicio se evaluó?\\
Distribución & ¿Cómo se asignan documentos a particiones o regiones?\\
Consistencia & ¿Qué garantía observa la consulta medida?\\
Índices & ¿Qué campos, orden y ámbito cubre cada índice?\\
Costo & ¿Qué operaciones, bytes y almacenamiento se facturan?\\
Recuperación & ¿Cómo se exportan, reconstruyen y verifican documentos?\\
\bottomrule\end{tabularx}\end{table}

\chapter{Modelo documental neutral}

\section{Documento canónico}

El documento de referencia contiene identidad estable, ámbito, revisión, versión de
esquema, tiempo de validez, tiempo de registro, estado y carga variable. Los campos
opcionales se distribuyen según tasas congeladas. Las claves no incluyen información
personal y se derivan de la semilla del generador.

\section{Proyección de cobertura}

La proyección conserva identidad, ámbito, revisión, campos de filtro, orden y salida de
la consulta objetivo. No contiene reglas de negocio ni acepta escrituras autoritativas.
Cada elemento enlaza documento fuente y versión de transformación. Una variante más
ancha se permite como análisis de sensibilidad, pero se publica como configuración
distinta.

\section{Catálogos y evolución}

Los datos incluyen versiones antiguas, valores desconocidos y sustituciones de catálogo.
La transformación debe migrar o aislar cada caso según una política declarada. El
generador conserva la verdad esperada para que la limpieza pueda evaluarse sin inspección
manual de todas las filas lógicas.

\chapter{Generación de cargas semirrealistas}

\section{Perfiles}

El perfil N1 representa eventos y multimedia; N2, reservas y disponibilidad; N3,
inventario, ventas y cierres. Cada perfil modifica selectividad, tamaño, frecuencia de
actualización, eventos tardíos y campos opcionales. Las diferencias deben afectar la
forma de la carga sin revelar nombres ni reglas de los sistemas que la inspiraron.

\section{Anomalías controladas}

\begin{table}[H]\centering
\caption{Anomalías que debe producir el generador.}
\begin{tabularx}{\textwidth}{p{3.4cm}p{3cm}X}\toprule
Anomalía & Parámetro & Resultado esperado\\\midrule
Campo ausente & tasa por perfil & valor predeterminado o exclusión declarada\\
Documento antiguo & versión de esquema & migración determinista\\
Evento tardío & retraso máximo & corte temporal reproducible\\
Entrega duplicada & identidad lógica & un solo efecto\\
Catálogo desconocido & tasa de novedad & cuarentena observable\\
Revisión obsoleta & distancia de revisión & rechazo sin regresión\\
\bottomrule\end{tabularx}\end{table}

La semilla gobierna identidades, orden y anomalías. Cambiarla produce otra muestra, no
otra definición del experimento. Toda tasa se registra en el manifiesto y se incluye en
el identificador de configuración.

\chapter{Protocolo P1: lectura documental}

\section{Configuraciones}

N0 consulta documentos canónicos usando el índice nativo apropiado. N1 consulta una
proyección de cobertura reconstruible. N2, opcional, evalúa una proyección más ancha para
estimar el costo de duplicar campos. No se utiliza una consulta deliberadamente sin
índice como adversario principal, porque la pregunta científica es si la proyección
mejora una implementación NoSQL razonable.

\section{Ejecución}

Cada corrida prepara un conjunto nuevo o restaura una instantánea congelada, verifica
conteos lógicos, calienta según protocolo y rota el orden de configuraciones. Se registran
latencias por consulta, documentos examinados cuando el motor lo permite, bytes de
respuesta, operaciones facturables y marca de frescura. La corrida se resume mediante
p50, p95 y p99.

\section{Interpretación}

Una diferencia pequeña de latencia puede ser irrelevante si la proyección duplica datos
sensibles, aumenta costo o incumple frescura. El informe presenta un vector de efectos.
La recomendación se formula por región de carga y nunca como superioridad universal.

\chapter{Protocolo P2: escritura y convergencia}

\section{Camino de actualización}

P2 modifica el documento canónico con revisión esperada. El mecanismo produce o detecta
el cambio, actualiza el derivado de forma idempotente y publica la nueva marca de
frescura. Se miden tiempo de confirmación canónica y tiempo de convergencia por separado.

\section{Amplificación}

La amplificación incluye documentos escritos, eventos, reintentos, bytes, invocaciones,
lecturas auxiliares, almacenamiento y tiempo operacional. Una implementación que oculta
trabajo en una función o canal de cambios no posee amplificación cero: debe atribuir ese
costo a la configuración.

\section{Frontera de decisión}

Sea $\Delta Q$ el ahorro por lectura y $\Delta W$ el costo adicional por escritura. La
frontera temporal simple exige $Q\Delta Q>W\Delta W$. El informe añade costo monetario,
almacenamiento, frescura y riesgo como dimensiones separadas. Si $\Delta Q\leq0$, la
proyección no se justifica por esa consulta.

\chapter{Protocolos de reconstrucción y publicación}

\section{Reconstrucción total}

R1 crea una versión candidata desde documentos canónicos, catálogos y transformación
congelada. La candidata no es visible para consumidores hasta superar conteo, hash por
partición, muestreo de contenido e invariantes de ámbito. La publicación cambia un alias
o registro de versión y conserva la anterior durante una ventana reversible.

\section{Reconstrucción incremental}

La ruta incremental consume cambios posteriores a un punto de control. Cada lote registra
inicio, fin, documentos procesados, duplicados y revisión máxima. Una reconstrucción
total periódica actúa como verificador independiente. Si ambas rutas divergen, la
incremental se considera defectuosa aunque la consulta parezca correcta.

\section{Interrupción segura}

El arnés interrumpe preparación, escritura parcial, verificación y cambio de versión. En
ningún caso una candidata incompleta debe reemplazar la versión activa. La reanudación
puede continuar o reiniciar, pero siempre produce el mismo resultado lógico.

\chapter{Aislamiento, autorización y recuperación}

\section{Matriz adversarial}

La prueba combina identidad válida, identidad ausente, ámbito manipulado, rol insuficiente,
clave de caché incompleta, tarea administrativa y recuperación semántica. Los candidatos
se restringen antes de ordenar o construir contexto. Un resultado ocultado después de
recuperarlo ya constituye una violación del contrato.

\section{Pruebas negativas}

Cada ruta positiva exige una negativa equivalente. Se verifican identificadores de otro
ámbito, referencias cruzadas, eventos reinyectados, paginación y búsqueda por campos
opcionales. Los registros de auditoría evitan contenido sensible, pero conservan motivo,
política, actor, ámbito y correlación.

\chapter{Productos para minería de datos e inteligencia artificial}

\section{Conjunto reproducible}

El conjunto declara población, corte temporal, fuente, versiones, limpieza, catálogos,
exclusiones, semilla y división. Su identificador cambia cuando una dependencia cambia.
La réplica reconstruye el conjunto desde cero y compara manifiesto, conteos, distribución
y huellas de contenido.

\section{Recuperación con evidencia}

Cada fragmento enlaza documento, revisión, ámbito, transformación y autorización. Se
miden precisión, exhaustividad, abstención, fugas y citas válidas. Las instrucciones
contenidas en documentos se tratan como datos, no como órdenes. La calidad lingüística
de una respuesta requiere un estudio separado.

\section{Herramientas de agentes}

Una herramienta declara capacidad, esquema de entrada, validación, idempotencia, efecto,
confirmación y auditoría. Las pruebas incluyen repetición, entrada parcial, cambio de
ámbito y cancelación. Una acción de alto riesgo sin confirmación falla el protocolo aun
si el resultado funcional parece correcto.

\chapter{Observabilidad y costos}

\section{Indicadores}

El tablero mínimo contiene latencia, error, frescura, atraso, duplicados, rechazos de
revisión, reconstrucciones, divergencia, operaciones por solicitud, almacenamiento y
costo por ámbito. Los indicadores se segmentan por versión de proyección para detectar
mezclas durante migración.

\section{Presupuesto de proyecciones}

Cada derivado consume un presupuesto compuesto por almacenamiento, escritura, eventos,
monitoreo y carga cognitiva. El responsable revisa uso y beneficio en una fecha fija.
Una proyección sin consultas, sin verificador o con reconstrucción rota entra en retiro.

\section{Criterio de retiro}

Retirar significa detener consumidores, verificar alternativas, eliminar el productor,
esperar la ventana de seguridad y suprimir datos derivados. La fuente canónica permanece
intacta. El proceso se prueba como cualquier migración y conserva un registro de decisión.

\chapter{Análisis estadístico y resultados negativos}

\section{Unidad inferencial}

La unidad es la corrida completa. Las consultas internas estiman un percentil de esa
corrida y no se presentan como observaciones independientes. Las comparaciones usan
pares generados con el mismo conjunto y semilla. Se informan mediana, intervalo
remuestreado, diferencia, razón y dirección de cada par.

\section{Multiplicidad y sensibilidad}

Las familias de hipótesis se congelan antes de medir y se ajustan por comparaciones
múltiples. El análisis de sensibilidad cambia una dimensión por vez: escala, tamaño de
documento, selectividad, consistencia o ancho de proyección. Las exploraciones se marcan
como tales y no reemplazan el análisis confirmatorio.

\section{Publicación de derrotas}

Una celda donde el documento canónico indexado iguala o supera a la proyección es un
resultado útil. Delimita dónde COBRIA debe simplificarse. El paquete conserva corridas
válidas aunque contradigan HN5 y explica cualquier descarte mediante una regla previa.

\chapter{Paquete abierto y revisión independiente}

\section{Contenido del depósito}

El depósito incluye protocolo, enmiendas, generador, adaptadores, manifiestos, eventos
crudos, análisis, tablas, figuras, entorno, licencias y declaración de disponibilidad.
Un comando de verificación comprueba huellas y vuelve a producir las salidas publicadas.

\section{Revisión técnica}

La persona revisora recibe una copia anonimizada, rúbrica y matriz de trazabilidad. Debe
examinar equivalencia de configuraciones, unidad inferencial, exclusiones, afirmaciones y
amenazas. Sus observaciones se responden una por una sin borrar desacuerdos sustantivos.

\section{Criterio de cierre}

La fase cuantitativa se considera completa cuando dos motores NoSQL terminan todas las
celdas congeladas, el paquete reproduce cifras y una revisión independiente no detecta
una falla que invalide la inferencia principal. Antes de ese punto, los resultados se
describen como piloto o evidencia parcial.

\chapter{Cuadernos de campo reproducibles}

\section{Bitácora de preparación}

Antes de ejecutar una celda, la bitácora registra identificador, fecha, responsable,
motor NoSQL, versión, configuración efectiva, semilla, conjunto, estado de índices y
huellas de los artefactos. El registro no depende de memoria personal: se produce desde
el mismo entorno que ejecuta la prueba y se conserva junto al resultado crudo. Si una
configuración cambia, nace una celda nueva; nunca se sobrescribe la evidencia anterior.

La preparación verifica además espacio disponible, sincronización temporal, ausencia de
carga ajena y aislamiento de credenciales. Una lectura breve de calentamiento confirma
que el adaptador puede resolver documento, ámbito y revisión. Este control evita gastar
una corrida completa para descubrir un error elemental de montaje.

\section{Bitácora de ejecución}

Cada fase deja marcas de inicio y final, conteos aceptados y rechazados, memoria máxima,
operaciones por segundo, percentiles de latencia y códigos de error. Los mensajes libres
se complementan con campos estructurados para permitir análisis automático. El operador
puede anotar un incidente, pero no clasificarlo retroactivamente sin conservar la razón y
la regla de exclusión utilizada.

\section{Bitácora de cierre}

Al terminar se verifica que los conteos coincidan con el manifiesto, que no existan
ámbitos cruzados, que las colas estén drenadas y que las huellas puedan recalcularse.
Una corrida válida se sella como inmutable. Una corrida inválida también se conserva,
etiquetada con su causa, porque la tasa y naturaleza de los fallos forman parte de la
evaluación de operabilidad.

\chapter{Migraciones y evolución del contrato}

\section{Cambio aditivo}

El cambio preferido agrega un campo opcional con semántica, valor ausente explícito y
versión de contrato. Lectores nuevos aceptan documentos antiguos; lectores antiguos
ignoran el campo nuevo cuando esa conducta es segura. La prueba mínima cubre ambas
direcciones y demuestra que la ausencia no se confunde con cero, cadena vacía o falso.

\section{Cambio incompatible}

Renombrar, dividir o reinterpretar un campo exige lectura doble, escritura controlada y
una ventana de convivencia. La migración se expresa como función determinista e
idempotente. Antes de retirar la versión previa se reconstruyen proyecciones, se comparan
resultados por ámbito y se comprueba que ningún consumidor siga solicitando el contrato
anterior.

\section{Reversión}

Toda migración declara el punto hasta el cual puede revertirse sin pérdida. Si el cambio
destruye información, se conserva una copia temporal gobernada y con caducidad. Revertir
no significa restaurar archivos a ciegas: significa volver a una combinación coherente
de contrato, documentos, catálogos, índices, proyecciones y consumidores.

\chapter{Rúbricas de conformidad COBRIA}

\section{Rúbrica del documento canónico}

La conformidad se valora de cero a dos por criterio: identidad estable, ámbito explícito,
revisión monotónica, procedencia, contrato versionado y validación en frontera. Cero
indica ausencia; uno, presencia parcial o no verificable; dos, evidencia automática. Un
promedio alto no compensa un cero en aislamiento o autoridad, que son condiciones de
seguridad y no preferencias de estilo.

\section{Rúbrica del derivado}

Un derivado demuestra propietario, propósito de consulta, versión de proyección,
idempotencia, frescura observable, reconstrucción total y procedimiento de retiro. Se
evalúa también si puede confundirse con la fuente autorizada. La etiqueta ``caché'' no
exime ningún criterio: si persiste y alimenta decisiones, debe poder explicarse y rehacerse.

\section{Rúbrica del uso analítico}

El conjunto analítico enlaza cada fila o fragmento con documento, revisión, ámbito y
transformación. La partición de entrenamiento evita contaminación temporal y cruce de
entidades. Una salida de inteligencia artificial conserva citas comprobables y se
abstiene cuando no reúne evidencia autorizada suficiente. La métrica de calidad se
publica junto con la de fuga, no en sustitución de ella.

\chapter{Escenarios de aceptación integral}

\section{Escenario A: lectura y actualización concurrentes}

Se inicia una lectura sobre una revisión conocida mientras llega una actualización del
mismo documento. El sistema puede devolver la revisión anterior o la nueva según el
contrato declarado, pero nunca una mezcla de campos. El evento posterior incluye la
nueva revisión, la proyección converge y una repetición no duplica efectos.

\section{Escenario B: reconstrucción bajo cambio de esquema}

Se elimina un derivado, se activa una nueva versión de proyección y se recorre la fuente
canónica. El resultado se compara con una construcción incremental limpia usando huellas,
conteos y consultas de aceptación. Toda diferencia queda localizada por documento y
transformación; una coincidencia meramente global no basta.

\section{Escenario C: recuperación con inteligencia artificial}

Una consulta solicita información disponible en dos ámbitos, pero la identidad solo
autoriza uno. El recuperador filtra antes de clasificar, conserva procedencia y entrega
únicamente fragmentos permitidos. Si esos fragmentos no sustentan una respuesta, el
sistema se abstiene. La prueba falla ante una cita inexistente, una revisión obsoleta no
declarada o cualquier exposición del ámbito prohibido.

\section{Escenario D: degradación y recuperación}

Se interrumpe temporalmente el consumidor que mantiene una proyección. La fuente acepta
escrituras dentro de su contrato, el atraso se vuelve visible y las lecturas degradadas
se etiquetan. Al reanudar, el consumidor procesa desde el último punto confirmado,
converge sin duplicados y demuestra equivalencia con una reconstrucción independiente.
